% =====================================================================
% El preprocesamiento del dato
% Parte I. El dato
%
% ENCARGO: el pipeline etapa por etapa. Es el capitulo central del
% documento: cada etapa con su figura de antes y despues.
%
% Reglas del documento:
% - Una idea = una subseccion = una figura.
% - Toda figura declara su procedencia con.
% - Toda cifra sale del canon de agosto; ver GUIA.md.
% - M1 y M3 se reportan siempre en paralelo.
% - Ningun superlativo sin releer el CSV de la frontera correcta.
% =====================================================================
\section{El preprocesamiento del dato}
\label{sec:prep}
\justifying

El capítulo anterior describió qué se midió y con qué instrumentos. Este
detalla lo que le ocurre a esas lecturas antes de que el modelo las
reciba.

La distancia entre los dos extremos es amplia: de un lado están los
archivos que escribe cada uno de los 27 equipos, con una lectura cada dos
minutos y huecos allí donde la adquisición falló. Del otro lado está lo
que el modelo necesita para operar, que es cuánto consumió y cuánto generó
cada institución en cada una de las \num{6144} horas del horizonte. Entre
ambos median seis etapas, y ninguna es automática: en todas se decidió
algo. Juzgar la solidez de las cifras del
Capítulo~\ref{sec:res} exige, por tanto, haber recorrido antes ese
proceso completo.

\subsection{Estructura del canal de preprocesamiento}
\label{sub:prep-mapa}

Antes de precisar cada etapa conviene ver el flujo completo de
procesamiento, porque las seis no son todas de la misma clase: tres
alteran el valor del dato y las otras tres lo preparan, lo ensamblan o lo
verifican. La
Figura~\ref{fig:prep-pipeline} las ordena y las distingue.

\begin{figure}[H]
 \centering
 \begin{tikzpicture}[
 x=1cm, y=1cm,
 font=\linespread{1}\scriptsize\selectfont,
 etapa/.style={anchor=west, draw=black!45, fill=white,
 rounded corners=2pt, minimum width=2.30cm, minimum height=1.12cm,
 align=center, inner sep=2pt},
 clave/.style={etapa, draw=infocolor, fill=infocolor!14, thick},
 salida/.style={etapa, draw=successcolor, fill=successcolor!12,
 thick, minimum width=4.96cm},
 chip/.style={anchor=west, rounded corners=2pt, thick,
 minimum width=1.10cm, minimum height=0.42cm, inner sep=1pt},
 flecha/.style={-{Stealth[length=4.5pt, width=3.5pt]}, draw=black!55,
 line width=0.5pt},
 retorno/.style={flecha, draw=black!40, rounded corners=3pt},
 rotulo/.style={anchor=east, align=right, text=black!62,
 font=\linespread{1}\scriptsize\selectfont\itshape},
 nota/.style={align=left, text=black!62,
 font=\linespread{1}\scriptsize\selectfont\itshape},
 ]
 % Tres carriles. El del medio reúne las etapas que alteran el valor del
 % dato; los otros dos, la que lo prepara y las que lo verifican. Se
 % pintan al fondo para que las cajas queden encima.
 \begin{scope}[on background layer]
 \fill[black!4, rounded corners=3pt] (2.80,-0.72) rectangle (6.55,0.72);
 \fill[infocolor!8, rounded corners=3pt] (2.80,-2.67) rectangle (16.05,-1.23);
 \fill[black!4, rounded corners=3pt] (2.80,-4.62) rectangle (16.05,-3.18);
 \end{scope}

 % --- Carril 1: del archivo a la serie horaria ------------------------
 \node[etapa, minimum width=3.05cm] (e1) at (2.9500, 0)
 {\textbf{1}\\Del archivo a\\la serie horaria};

 % --- Carril 2: las tres etapas que alteran el valor ------------------
 \node[clave] (e2) at (2.9500, -1.95) {\textbf{2}\\Clasificar\\el medidor};
 \node[clave] (e3) at (5.6125, -1.95) {\textbf{3}\\Reconstruir la\\demanda bruta};
 \node[clave] (e4) at (8.2750, -1.95) {\textbf{4}\\Atípicos\\y huecos};

 % --- Carril 3: ensamblar, verificar y entregar -----------------------
 \node[etapa] (e5) at (2.9500, -3.90) {\textbf{5}\\Ensamblar\\y verificar};
 \node[etapa] (e6) at (5.6125, -3.90) {\textbf{6}\\Fijar la\\zona horaria};
 \node[salida] (out) at (10.9375,-3.90)
 {\textbf{Matrices $D$ y $G$}\\$5 \times \num{6144}$, no negativas};

 % --- Rótulos de carril ----------------------------------------------
 \node[rotulo] at (2.62, 0) {transporta,\\ordena\\y convierte};
 \node[rotulo, text=infocolor,
 font=\linespread{1}\scriptsize\selectfont\bfseries]
 at (2.62,-1.95) {las tres etapas\\que alteran\\el valor del dato};
 \node[rotulo] at (2.62,-3.90) {ensamblan\\y verifican};

 % --- Avance dentro de cada carril ------------------------------------
 \draw[flecha] (e2) -- (e3);
 \draw[flecha] (e3) -- (e4);
 \draw[flecha] (e5) -- (e6);
 \draw[flecha] (e6) -- (out);

 % --- Paso de un carril al siguiente ----------------------------------
 \draw[retorno] (e1.south) -- (e2.north);
 \draw[retorno] (e4.south) -- ++(0,-0.42) -| (e5.north);

 % --- La bifurcación de las dos fronteras, en la etapa 1 --------------
 \node[chip, draw=sectioncolor, fill=sectioncolor!14, text=sectioncolor]
 (pm1) at (2.95, 1.34) {\textbf{M1}};
 \node[chip, draw=warningcolor, fill=warningcolor!16, text=warningcolor]
 (pm3) at (4.15, 1.34) {\textbf{M3}};
 \draw[flecha, draw=sectioncolor, line width=0.7pt]
 (pm1.south) -- (e1.north);
 \draw[flecha, draw=warningcolor, line width=0.7pt]
 (pm3.south) -- (e1.north);
 \node[nota, anchor=west, text width=8.4cm] at (6.75, 1.34)
 {aquí se separan las dos fronteras de medición:
 cambia el medidor que se lee y, con él, su tipo};

 \node[nota, anchor=center, align=center, text=infocolor]
 at (13.31,-1.95) {aquí se juega\\la fidelidad del dato};
 \end{tikzpicture}
 \caption[El pipeline de preprocesamiento en seis etapas]{Las 6
 etapas del preprocesamiento, en el orden en que se ejecutan. La
 bifurcación entre las dos fronteras de medición ocurre en la etapa~1.}
 \label{fig:prep-pipeline}
\end{figure}

Esas tres se reparten a su vez en dos grupos: la primera lleva el archivo
hasta la serie horaria, y las dos últimas la ensamblan y la verifican. Las
tres resaltadas en la figura son las que modifican el valor. Esa separación es
la que permite decir con precisión dónde se juega la fidelidad del dato.

Del mapa importa retener una propiedad. Todo el trabajo se recorre dos
veces, una por cada frontera de medición, y lo único que las distingue es
qué medidor se lee en la primera etapa. De ahí en adelante el tratamiento es
el mismo para las dos: si algo sale distinto es porque el medidor elegido
lo provoca, no porque se le aplique una regla diferente.

Cualquier diferencia entre M1 y M3 en los resultados es, por tanto,
atribuible a qué se mide y no a cómo se procesa. Hay una excepción, y
es importante definirla ahora: Mariana no tiene un segundo medidor
que leer, de modo que en la frontera secundaria se reutiliza el primero
multiplicado por un factor que el Capítulo~\ref{sec:frontera} declara.

Las seis etapas no se reparten por igual en lo que sigue. La primera ocupa
cuatro subsecciones, porque es donde el dato cambia de paso y de forma. La
segunda y la tercera ocupan dos cada una, y en los dos casos la primera de
ellas está dedicada al fenómeno que las motiva. Las tres últimas ocupan una
cada una. A partir de ahí el capítulo deja de describir etapas: las
subsecciones finales leen el resultado y declaran qué le hacen a la cifra
las decisiones tomadas por el camino.

\subsection{Del archivo a la serie horaria}
\label{sub:prep-lectura}

La primera etapa no interviene sobre el valor medido: lo transporta, lo
ordena y lo expresa en el paso y las unidades con que el modelo trabaja.
Recibe los archivos de un equipo, con una lectura cada dos minutos, y
entrega una serie de una hora. La Figura~\ref{fig:prep-archivo-serie}
recorre ese camino sobre una hora real de la UCC, es decir, sigue las
lecturas del medidor y las del inversor desde el archivo en que llegaron
hasta el valor horario que cada una deja en la serie.

\begin{figure}[H]
 \centering
 \includegraphics[width=0.95\textwidth]{figuras/f3_01_archivo_a_serie_m1.png}
 \caption[Del archivo a la serie horaria, sobre una hora real]{El
 recorrido de una hora real, la de las 13 horas del jueves 6 de noviembre
 de 2025 en la UCC, frontera M1.}
 \label{fig:prep-archivo-serie}
\end{figure}

Los dos equipos no hablan en la misma unidad. El medidor registra en
kilovatios y el inversor en vatios, de modo que sus series no se pueden
comparar ni sumar tal como llegan. La del inversor se divide entre mil y
la del medidor se usa como viene.

La etapa lleva además cuatro salvaguardas contra cosas que este
conjunto no contiene, y de ellas lo que importa es cuántas veces actúan:
ninguna.

\begin{enumerate}[leftmargin=*, itemsep=0.35em]
 \item \textbf{Archivos que no se dejan leer.} No hay ninguno entre los 73
 del árbol.

 \item \textbf{Fechas ilegibles.} No hay ninguna entre los \num{2813305}
 registros de los diez medidores.

 \item \textbf{Valores no numéricos.} Tampoco hay ninguno entre esos mismos
 registros.

 \item \textbf{Lecturas negativas del inversor.} No hay ninguna entre las
 \num{793981} de los siete equipos.
\end{enumerate}

La última no actúa porque un inversor entrega energía a la red y no la toma
de ella. Los siete registran el cero miles de veces, de modo que la
ausencia de valores negativos no se debe a que el equipo no llegue hasta
ahí, sino a que por debajo del cero no hay nada que medir.

Queda una sola ambigüedad que la etapa sí tiene que resolver, y es la del
instante que llega dos veces.

\subsection{Los instantes duplicados}
\label{sub:prep-duplicados}

Antes de promediar hay que resolver una ambigüedad del dato crudo. En la totalidad de la base de datos algunas
lecturas comparten instante con otra, \num{3108} veces en la frontera
principal y \num{3588} en la secundaria, pero el registro de un equipo no
puede llevar dos valores en el mismo instante. Esto se resuelve dejando uno solo,
tomando el promedio entre sí de los datos duplicados. En el horizonte de estudio solo quedan 27 en cada frontera.

El fenómeno es de un equipo y no del conjunto: \num{3073} de los
\num{3108} instantes de la frontera principal son del medidor de la UCC, y
en las otras cuatro instituciones no pasan de 30 apariciones en todo el
archivo. Por eso el ejemplo de la Figura~\ref{fig:prep-archivo-serie} es de
la UCC, que es la única institución donde la fusión llega a ocurrir dentro
del horizonte.

La regla para determinar qué instante de los datos duplicados usar no tiene consecuencia. En los \num{6696} instantes repetidos de las dos fronteras las
lecturas que lo comparten coinciden hasta el último dígito, de modo que
promediarlas, quedarse con la primera o quedarse con la última dan
exactamente el mismo número. Se toma el promedio porque el instante ha de ser único para poder alinear los tramos
de un mismo medidor, no porque haya dos valores en disputa que arbitrar.

La regla vale igual dentro de un archivo y entre archivos. Los tramos de un
mismo medidor comparten un instante en su costura, una vez por institución
y por frontera, y también ahí la lectura repetida se funde en una sola.

\begin{figure}[H]
 \centering
 \includegraphics[width=0.95\textwidth]{figuras/f3_01d_duplicados.png}
 \caption[Cuándo y en qué equipo se repiten los instantes]{Los instantes
 repetidos de las dos fronteras de medición, en el tiempo y por
 institución.}
 \label{fig:prep-duplicados}
\end{figure}

\subsection{La energía de la hora}
\label{sub:prep-energia}

La serie se remuestrea entonces a una hora tomando la media de sus
muestras, y esa media no es una elección de método: sale de la definición
de energía. La energía de la hora es la integral de la potencia sobre ella,
según la ecuación~\eqref{eq:dato-energia}, y de esa integral solo se
conocen valores separados por un paso fijo, de modo que el tramo observado
se suma por rectángulos de ese ancho:

\begin{equation}
 E \;=\; \int_{t}^{t+\Delta t} p(\tau)\,\mathrm{d}\tau
   \;\simeq\; \delta \sum_{k=1}^{n} P_{k}
   \;=\; \overline{P}\,(n\,\delta) ,
 \qquad \delta = \uni{2}{min},
 \label{eq:prep-media}
\end{equation}

donde $n$ es el número de muestras que la hora trae, $\delta$ el paso de
muestreo y $\overline{P}$ la media. Que el paso sea fijo no es un supuesto
sino un hecho medido: sobre \num{1117638} registros de cuatro medidores
vale exactamente dos minutos en más del
\pct{99.6} de los intervalos. De ahí que sumar los tramos y promediar las
muestras no sean dos métodos entre los que haya que elegir, sino la misma
operación escrita de dos maneras.

Con la hora completa $n$ vale 30, el producto $n\,\delta$ vale una hora, y
la energía en kilovatios hora resulta numéricamente igual a la potencia
media en kilovatios, de modo que la misma serie puede leerse de las dos
maneras y a lo largo del documento se usan ambas según convenga. La
Figura~\ref{fig:prep-energia} lo enseña sobre la hora del ejemplo.

\begin{figure}[H]
 \centering
 \includegraphics[width=0.95\textwidth]{figuras/f3_01b_energia_hora.png}
 \caption[La energía de una hora, contada de dos maneras]{Las treinta
 lecturas de la hora del ejemplo dibujadas como rectángulos de dos minutos,
 y encima el rectángulo único de la hora entera a la altura de la media. Las
 dos áreas son la misma área y los dos números coinciden hasta el último
 dígito, porque el paso de muestreo es constante.}
 \label{fig:prep-energia}
\end{figure}

\subsection{La hora incompleta}
\label{sub:prep-incompleta}

Cuando la hora llega incompleta, en cambio, los rectángulos no cubren el
intervalo y el tramo sin muestrear no queda determinado por el dato. Ese es
el único punto de la etapa en el que hay que suponer algo, y aquí se le
atribuye al tramo la potencia media de las muestras presentes, es decir, se
toma $E \simeq \overline{P}\,\Delta t$ con $\Delta t = \uni{1}{h}$. Es un
supuesto y no una identidad, de modo que se declara como tal y se acota en
la subsección~\ref{sub:prep-sesgo}.

Una hora puede llegar de tres maneras, y el tratamiento cambia en cada una.

\begin{enumerate}[leftmargin=*, itemsep=0.35em]
 \item \textbf{Completa}, con sus 30 muestras. La energía queda
 determinada por el dato y no hay nada que decidir.

 \item \textbf{Incompleta}, con entre 1 y 29 muestras. Se estima con la
 media de las presentes, que es el supuesto que se acaba de declarar. Son
 \num{835} horas del horizonte.

 \item \textbf{Vacía}, sin ninguna muestra. No se estima aquí: pasa a la
 etapa de limpieza, que interpola sus primeras horas y repite un valor en
 las restantes, según describe la subsección~\ref{sub:prep-huecos}. Son
 \num{567} horas.
\end{enumerate}

La media horaria ignora las muestras que faltan, de modo que la
estimación es lo que la agregación hace por defecto. La otra opción es contar solo los minutos medidos y no añadir nada por los
que faltan. Parece la opción prudente, la que no supone nada, pero supone
tanto como la anterior: da por hecho que en los minutos sin registro no se
consumió energía, y en instalaciones con carga de base permanente eso es
falso. La Figura~\ref{fig:prep-incompleta} pone las dos suposiciones frente
al contador de energía del propio medidor, que es el único dato capaz de
arbitrar entre ellas.

\begin{figure}[H]
 \centering
 \includegraphics[width=0.95\textwidth]{figuras/f3_01c_hora_incompleta.png}
 \caption[La hora incompleta y sus dos suposiciones]{A la izquierda, una
 hora que trae 18 de sus 30 muestras: suponer la media para el tramo que
 falta deja la hora en \uni{8.98}{kWh} y el contador del equipo marcó
 \uni{9.00}{kWh}, mientras que quedarse con lo observado la dejaría en
 \uni{5.39}{kWh}. A la derecha, la misma comparación sobre todas las horas
 incompletas de los tres medidores cuyo contador está bien escalado.}
 \label{fig:prep-incompleta}
\end{figure}

El regulador, ante el mismo problema, tampoco asigna cero. La Resolución
CREG 038 de 2014 ordena estimar la lectura mientras el sistema de medición
está en falla, y en las fronteras que no reportan al administrador del
mercado se estima por consumos promedios del mismo usuario, que es el
criterio que aquí se aplica. Rige sobre fronteras comerciales y no sobre
instrumentación de investigación como esta, de modo que se toma como
criterio y no como parte de la alineación con la normativa. El articulado y su alcance exacto quedan en el Anexo~\ref{anx:repro}.que recoge el procedimiento completo de la etapa, con esas cuatro comprobaciones y el orden exacto en que se aplican.

\begin{guardado}
\begin{detallebox}

Las dos decisiones se pueden contrastar, porque el medidor lleva junto a la
potencia sus propios contadores de energía importada y exportada, en
kilovatios hora, que avanzan por su cuenta y no intervienen en el
preprocesamiento. El contraste usa la diferencia entre los dos y no la
importada sola, por la razón que da la nota de la
Figura~\ref{fig:prep-incompleta}. Sobre las \num{5810} horas completas del
circuito principal de la UCC en las que el contador tiene lectura en los dos
bordes, lo que avanza y la media horaria difieren en el \pct{0.04}
acumulado.

El supuesto sobre el tramo no observado admite una prueba más directa, y es
la que decide. En los \num{718} cortes de telemetría de los tres medidores
cuyo contador está bien escalado, el contador avanza \uni{809}{kWh};
suponer que el equipo siguió consumiendo predice \uni{812}{kWh} y suponer
que no hubo energía predice \uni{288}{kWh}. El equipo sigue integrando
durante los cortes, de modo que atribuirles energía nula falla por física y
no por estadística, y falla cada vez peor cuanto más largo es el corte.

Hay además dos cantidades que se parecen y difieren por un factor de ocho. Las \num{835} horas incompletas del horizonte
\emph{contienen} \uni{8100}{kWh}, el \pct{3.2} de la demanda, pero de esa
energía casi toda está observada: lo que el supuesto \emph{imputa} es solo
el tramo que falta, y son \uni{1054}{kWh}, el \pct{0.41}. La razón es que
esas horas pierden muy poco. El \pct{55.6} pierde una sola muestra de las
treinta y el \pct{77} pierde cuatro o menos.

Queda por saber si la estimación descansa sobre las horas peor cubiertas, y
no lo hace. Tomando como válida solo la hora que trae al menos el \pct{75}
de sus muestras, \num{147} horas del horizonte quedarían fuera, es decir el
\pct{0.49} de las que tienen dato, y tratarlas como huecos en vez de
estimarlas mueve la demanda de la comunidad en \uni{65.6}{kWh}, el
\pct{0.026}. Ninguna institución se mueve más del \pct{0.42}, y esa es
Udenar.
\end{detallebox}
\end{guardado}



\subsection{La demanda negativa}
\label{sub:prep-negativa}

La primera etapa, la que lleva del archivo a la serie horaria, entrega
datos que no concuerdan: la demanda de tres de las cinco instituciones baja
de cero durante varias horas al día. Un edificio no consume energía
negativa. Y el tratamiento no lo causó, porque esa etapa no corrige, no
sustituye ni descarta ninguna medida, de modo que lo que la serie trae es
lo que el medidor entregó. La Figura~\ref{fig:prep-negativa} recoge el
fenómeno en sus dos dimensiones y en las dos fronteras: la forma que toma
en la institución más afectada de cada una y la hora del día en que
aparece.

\begin{figure}[H]
 \centering
 \includegraphics[width=0.95\textwidth]{figuras/f3_02_demanda_negativa.png}
 \caption[La demanda que el medidor entrega en negativo]{Una fila por
 frontera de medición. A la izquierda, el perfil de la institución que
 carga con el fenómeno en cada una, en promedio y en mínimo de cada hora,
 sobre las otras cuatro en promedio y en gris; la banda sombreada es la
 zona imposible. A la derecha, cuántas horas caen bajo cero en cada hora
 del día, apiladas por institución.}
 \label{fig:prep-negativa}
\end{figure}

\notafig{Las dos curvas de la institución destacada van juntas a
propósito. Bajo M1 el promedio horario de Udenar entra en la zona
imposible y baja a \uni{-11.0}{kW} a las 13, de modo que el neteo no
arrastra unos pocos casos sueltos sino la media entera; bajo M3 ningún
promedio la alcanza y solo el mínimo la roza. Los conteos por institución,
la fracción, el mínimo y la energía hacia la red los da la
Tabla~\ref{tab:prep-tipos}. Elaboración propia desde el caché de estados
intermedios.}

El reparto es muy desigual. En Udenar la demanda cae bajo cero en cerca de
una cuarta parte del horizonte, es decir, el \pct{24.7} de las \num{6144}
horas y el \pct{25.1} de las que traen dato. En Mariana y en la UCC el
fenómeno es ocasional, y en el HUDN y en CESMAG no aparece ni una hora, de
modo que una sola institución concentra el \pct{83} de todas las horas
afectadas. La Tabla~\ref{tab:prep-tipos} reúne el conteo y el mínimo de
cada una.

Lo decisivo, sin embargo, no es cuántas horas son sino \emph{a qué hora del
día} ocurren. Si las lecturas negativas fueran ruido del sensor se
repartirían de manera aproximadamente uniforme a lo largo del día. Lo que
se cuenta hora a hora es lo contrario: las negativas dibujan la campana de
la producción solar, con su máximo de \num{288} horas a las 12 y \num{282}
a las 13, y no hay ni una sola entre las 18 y las 6.

Ese patrón identifica la causa sin necesidad de abrir el equipo. No es un
sensor averiado, es un medidor que está restando la generación antes de
reportar, y dónde esté puesto respecto de la inyección es lo que decide si
lo hace.

Y decide también en qué frontera aparece, que es lo que separa las dos
filas de la figura. En la secundaria el fenómeno casi desaparece: de las
cinco instituciones solo Mariana registra lecturas negativas, y son las
mismas \num{213} horas de la principal, porque en M3 se lee su mismo
medidor multiplicado por un factor. Los otros cuatro circuitos secundarios
no bajan de cero ni una hora, de modo que la lectura imposible es un
asunto de la frontera principal.

\subsection{Clasificación del medidor: tres tipos}
\label{sub:prep-tipos}

La explicación de ese patrón es de topología eléctrica, o dicho de otro
modo, de en qué punto del circuito se conecta cada inversor. Cuando el
inversor entra \emph{entre el medidor y la carga}, la energía que produce
no llega a pasar por el medidor, de manera que este no registra la demanda
del circuito sino lo que queda de ella al descontarle la generación, es
decir, el flujo neto $D - G$. En las horas en que la generación supera la
carga del circuito, ese flujo se invierte de sentido y el medidor registra
un valor negativo. La lectura no es un
error: es la evidencia de un autoconsumo que el medidor no puede separar.

Según dónde entre el inversor, el tratamiento se bifurca en dos: a unos
medidores hay que devolverles la generación que descontaron y a otros no
hay nada que devolverles. El documento nombra además un caso intermedio,
el neto parcial, pero no porque reciba un tratamiento propio, que es el
mismo del neto, sino porque el neteo apenas llega a invertir el flujo. Lo
que separa a esos dos es el impacto.

\begin{figure}[H]
 \centering
 \begin{tikzpicture}[
 font=\scriptsize,
 eq/.style={draw=black!55, fill=black!3, rounded corners=1.5pt,
 minimum width=1.5cm, minimum height=0.58cm, align=center},
 med/.style={draw=warningcolor, fill=warningcolor!14, very thick,
 circle, minimum size=0.72cm, inner sep=0pt},
 inv/.style={draw=successcolor, fill=successcolor!14, very thick,
 rounded corners=1.5pt, minimum width=1.15cm,
 minimum height=0.55cm, align=center},
 lin/.style={-{Stealth[length=3.5pt]}, draw=black!60},
 et/.style={font=\scriptsize\itshape, text=black!65},
 ]
 % --- net ---
 \node[et] at (-2.35,0.62) {\textbf{neto}};
 \node[eq] (r1) at (-1.0,0) {Red};
 \node[med] (m1) at (0.55,0) {M};
 \node[eq] (c1) at (2.35,0) {Carga};
 \node[inv] (i1) at (2.35,-1.0) {Inversor};
 \draw[lin] (r1) -- (m1);
 \draw[lin] (m1) -- (c1);
 \draw[lin] (i1) -- (c1);
 \node[et, text width=3.4cm, right=0.30cm of c1, align=left]
 {el inversor entra entre\\el medidor y la carga:\\el medidor ve $D-G$};

 % --- gross ---
 \node[et] at (-2.35,-2.38) {\textbf{bruto}};
 \node[eq] (r2) at (-1.0,-3.0) {Red};
 \node[inv] (i2) at (0.55,-4.0) {Inversor};
 \node[med] (m2) at (2.05,-3.0) {M};
 \node[eq] (c2) at (3.6,-3.0) {Carga};
 \draw[lin] (r2) -- (m2);
 \draw[lin] (m2) -- (c2);
 \draw[lin] (i2) -- (r2);
 \node[et, text width=3.4cm, right=0.30cm of c2, align=left]
 {el inversor entra antes\\del medidor, que ve $D$\\sin descontar nada};
 \end{tikzpicture}
 \caption[Por qué unos medidores netean y otros no]{Posición relativa
 del medidor y del inversor en los dos casos extremos. Arriba, el inversor
 entra entre el medidor y la carga, de modo que el medidor observa el
 flujo neto, que se invierte en horas de alta irradiancia. Abajo, entra
 del lado de la red y el medidor registra la demanda bruta. El
 caso intermedio, el neto parcial, recibe el mismo tratamiento que el
 neto y se distingue solo por el grado, porque la generación que queda
 detrás del medidor pesa poco frente a la carga del circuito.}
 \label{fig:prep-tipos}
\end{figure}

\begin{table}[H]
 \centering
 \caption{Clasificación de los medidores de demanda y huella del neteo, en
 las dos fronteras. Las cuatro columnas de la derecha se leen directamente
 sobre la serie que cada medidor entrega, sin reconstruir nada: son horas
 en que su lectura fue negativa, la fracción que representan sobre las
 horas con dato de ese equipo, su lectura más baja del horizonte y la
 energía que en esas horas salió del circuito hacia la red.}
 \label{tab:prep-tipos}
 \small
 \begin{tabular}{@{}l l r r r r@{}}
 \toprule
 \textbf{Institución} & \textbf{Tipo} & \textbf{Horas} &
 \textbf{Sobre sus} & \textbf{Lectura} & \textbf{Energía hacia} \\
 & & \textbf{bajo cero} & \textbf{horas} & \textbf{mínima (kW)} &
 \textbf{la red (kWh)} \\
 \midrule
 \multicolumn{6}{@{}l}{\itshape Frontera M1} \\
 Udenar & neto & \num{1517} & \pct{25.1} & \num{-33,57} & \num{15348} \\
 Mariana & neto parcial & \num{213} & \pct{3.5} & \num{-2,41} & \num{156} \\
 UCC & neto parcial & \num{94} & \pct{1.6} & \num{-5,91} & \num{200} \\
 HUDN & bruto & \num{0} & \pct{0} & \num{6,12} & \num{0} \\
 CESMAG & bruto & \num{0} & \pct{0} & \num{0,20} & \num{0} \\
 \addlinespace
 \multicolumn{6}{@{}l}{\itshape Frontera M3} \\
 Udenar & bruto & \num{0} & \pct{0} & \num{0,22} & \num{0} \\
 Mariana & bruto & \num{213} & \pct{3.5} & \num{-0,72} & \num{47} \\
 UCC & bruto & \num{0} & \pct{0} & \num{0,73} & \num{0} \\
 HUDN & bruto & \num{0} & \pct{0} & \num{1,29} & \num{0} \\
 CESMAG & bruto & \num{0} & \pct{0} & \num{0,20} & \num{0} \\
 \bottomrule
 \end{tabular}
\end{table}

El bloque de M3 dice algo que no se deduce del de M1: en el circuito
secundario los cinco medidores son brutos, de modo que la frontera no
cambia solo cuánta carga se mide sino la semántica del propio equipo. La
única excepción aparente es Mariana, y no es tal: esa institución no tiene
medidor en el circuito secundario, de modo que se reutiliza el del
principal multiplicado por un factor, y con él hereda sus \num{213} horas
negativas, ya escaladas.

Los tres tipos no nombran tres aparatos distintos sino tres situaciones
del mismo, y lo que las separa es dónde se conecta la generación. Si entra
entre el medidor y la carga, el medidor no la ve llegar y registra solo la
demanda que queda después de restarla, de modo que su lectura puede caer
bajo cero. Si entra del lado de la red, el medidor la ve pasar y su lectura
nunca se invierte. La Figura~\ref{fig:prep-tipos} dibuja los dos casos.

La tabla ordena las cinco instituciones por lo que esa posición produce, y
la gradación se lee sin más. En Udenar el flujo se invierte una de cada
cuatro horas y llega a \uni{-33.57}{kW}. En Mariana y en la UCC ocurre en
el \pct{3.5} y el \pct{1.6} de sus horas, con mínimos veinte veces menos
profundos, de modo que la inversión es allí ocasional. En el HUDN y en
CESMAG no ocurre nunca, y la lectura más baja de CESMAG, \uni{0.20}{kW},
dice que su circuito llega a rozar el cero sin cruzarlo.

La columna de la derecha es la que mide el fenómeno sin ambigüedad, porque
no depende de ninguna reconstrucción: en el horizonte, del circuito de
Udenar salieron hacia la red \uni{15348}{kWh}, contra los \uni{156}{} de
Mariana y los \uni{200}{kWh} de la UCC.

En el HUDN y en CESMAG el inversor entra fuera del tramo que el medidor
vigila, o alimenta un circuito distinto del medido, y por eso el flujo
nunca se invierte allí aunque las dos instituciones tengan su equipo. La
Figura~\ref{fig:prep-profundidad} pone esa gradación en una sola magnitud,
es decir, sitúa a cada medidor entre la lectura que entrega de ordinario y
la más baja que llegó a entregar, y lo hace en las dos fronteras.

\begin{figure}[H]
 \centering
 \includegraphics[width=0.95\textwidth]{figuras/f3_02b_profundidad.png}
 \caption[Hasta dónde baja la lectura de cada medidor]{Una fila por
 institución en el orden fijo y un panel por frontera, con los kilovatios
 en el eje horizontal. Cada fila lleva la lectura mediana del medidor, el
 recorrido desde esa mediana hasta la más baja, y la mitad central de sus
 lecturas negativas cuando las tiene.}
 \label{fig:prep-profundidad}
\end{figure}

En Udenar esa mitad central cae entre
\uni{-15.1}{} y \uni{-3.9}{kW}, con mediana de \uni{-8.7}{kW}, de modo que
el mínimo de \uni{-33.57}{kW} casi cuadruplica la inversión ordinaria; en
Mariana la mediana de las negativas es \uni{-0.59}{kW} y en la UCC,
\uni{-1.94}{kW}. Cuántas horas resume cada fila lo da la
Tabla~\ref{tab:prep-tipos}.

En la frontera M3 la clasificación no hace falta, porque allí los cinco
medidores son brutos y ninguna institución requiere reconstrucción. La
excepción es Mariana, y no porque su circuito secundario netee: es que no
tiene un segundo medidor, de modo que su serie de M3 se construye
multiplicando la del primero por un factor, conforme declara el
Capítulo~\ref{sec:frontera}. Con la serie viajan sus \num{213} horas
negativas, reducidas por ese mismo factor, y llevarlas a cero cuesta
\uni{46.7}{kWh}, el \pct{0.33} de su demanda en esa frontera.

\subsection{La reconstrucción de la demanda bruta}
\label{sub:prep-reconstruccion}

Identificada la causa, la corrección es sencilla de enunciar: si el medidor
restó la generación, se le devuelve. El modelo necesita la demanda bruta
$D$ (lo que el circuito medido consume) y la generación $G$ por separado, de
modo que para los medidores netos y los netos parciales la demanda se
reconstruye como

\begin{equation}
 D_{\text{recon}}
 = \Bigl(D_{\text{neto}} + \sum_{\text{inversores}} G \Bigr)^{+},
 \label{eq:prep-reconstruccion}
\end{equation}

donde la suma recorre \emph{todos} los inversores físicos que inyectan
entre ese medidor y su carga, y el exponente $+$ significa que, si el
resultado sale negativo, se toma cero.

La expresión no se aplica por igual a las cinco instituciones. La
Tabla~\ref{tab:prep-reconstruccion} reparte los cinco casos: a quién se le
devuelve generación, con cuántos inversores y si el recorte a cero llega a
actuar.

\begin{table}[H]
 \centering
 \caption{Cómo se aplica la reconstrucción en cada institución. El HUDN y
 CESMAG entregan lectura bruta, de modo que no hay generación que
 devolverles y el recorte queda solo como resguardo. Cuando una hora no
 trae lectura de medidor pero sí de inversor, la demanda de partida se
 toma como nula y el resultado queda igual a la generación devuelta; pasa
 en \num{7} horas del horizonte, todas de Udenar, y suma \uni{34.1}{kWh}.}
 \label{tab:prep-reconstruccion}
 \small
 \begin{tabular}{@{}l l r r r@{}}
 \toprule
 \textbf{Institución} & \textbf{Tipo} & \textbf{Inversores} &
 \textbf{Horas con} & \textbf{Energía} \\
 & & \textbf{devueltos} & \textbf{recorte} &
 \textbf{recortada (kWh)} \\
 \midrule
 Udenar & neto & \num{3} & \num{353} & \num{1013,3} \\
 \addlinespace
 Mariana & neto parcial & \num{1} & \num{6} & \num{0,6} \\
 \addlinespace
 UCC & neto parcial & \num{1} & \num{0} & \num{0} \\
 \addlinespace
 HUDN & bruto & --- & \num{0} & \num{0} \\
 \addlinespace
 CESMAG & bruto & --- & \num{0} & \num{0} \\
 \bottomrule
 \end{tabular}
\end{table}

La única fila que pide explicación es la primera, porque los tres equipos
de Udenar no cumplen el mismo papel. El inversor del proyecto entró en
servicio el 3 de septiembre de 2025 y cubre \pct{22.9} del horizonte; los
otros dos registran desde el primer día. Antes de esa fecha, entonces, la
suma solo puede devolver lo que midieron esos dos, y queda incompleta.

Que sumar los tres sea lo correcto lo dice el propio medidor. Desde el 3
de septiembre, devolverle la generación de los tres deja
\emph{exactamente} cero horas con demanda negativa. Las \num{353} de la
tabla caen todas antes, que es donde la suma va incompleta.

Tomar como nula la demanda de partida tiene además una consecuencia de
mucho más alcance, y es que convierte en cero toda hora sin lectura de
medidor. Bajo M1 eso alcanza a \num{330} horas: \num{90} de Udenar,
\num{114} de Mariana y \num{126} de la UCC. Entran al modelo como consumo
nulo sin pasar por la etapa de limpieza, que ya no las ve, y sin que
ninguna marca las distinga de una hora realmente medida.

La Figura~\ref{fig:prep-reconstruccion} confronta las tres series sobre un
mismo eje y muestra el resultado de aplicar esa devolución. La coincidencia
temporal es el argumento: la lectura se hunde exactamente cuando la banda
de generación se ensancha, y ninguna avería de sensor produciría eso.

La Figura~\ref{fig:prep-mosaico} repite esa lectura en las diez series, de
modo que se vea a quién se le aplica la devolución y a quién no.

\begin{figure}[H]
 \centering
 \includegraphics[width=0.95\textwidth]{figuras/f3_03b_reconstruccion_mosaico.png}
 \caption[La reconstrucción en las diez series]{Perfil medio horario de la
 lectura del medidor y de la demanda reconstruida en las cinco
 instituciones y las dos fronteras, con la generación devuelta como banda
 entre las dos curvas. Las escalas verticales son independientes por
 panel.}
 \label{fig:prep-mosaico}
\end{figure}

\notafig{Siete de los diez paneles salen sin banda porque su medidor no
netea y no hay generación que devolverles, de modo que la fila de la
frontera secundaria funciona como control de la principal. La excepción es
la Universidad Mariana bajo M3, donde tampoco hay devolución y sin embargo
el recorte a cero todavía quita \uni{46.7}{kWh}, que son los
\uni{156}{kWh} de la frontera principal multiplicados por el mismo factor
con que se construye su serie secundaria. Elaboración propia desde el
caché de estados intermedios.}

\begin{figure}[H]
 \centering
 \includegraphics[width=0.95\textwidth]{figuras/f3_03_reconstruccion_m1.png}
 \caption[La reconstrucción de la demanda bruta de Udenar]{La
 reconstrucción sobre Udenar. A la izquierda, un viernes hora a hora: en
 rojo la lectura del medidor y en azul la demanda que resulta al
 devolverle la generación.}
 \label{fig:prep-reconstruccion}
\end{figure}

La banda beige entre las dos curvas es esa
generación, de modo que la curva azul es la roja levantada por la banda,
y la cota vertical la mide en la hora de lectura más baja, cuando los
tres inversores entregaron \uni{35.1}{kW}. El día pertenece al tramo
posterior al 3 de septiembre de \num{2025}, cuando los tres registran a
la vez, y ninguna de sus horas llega al recorte. A la derecha, el
promedio de las \num{6144} horas en la misma escala vertical. Abajo, la
misma igualdad en energía sobre el horizonte completo.

\begin{guardado}
\notafig{La curva azul resultante recupera la forma que cabe esperar de un
campus universitario, es decir, base nocturna estable, ascenso matinal,
valle de almuerzo y tarde sostenida. La franja inferior separa las dos
partes de la cuenta: de los \uni{33705}{kWh} en que la reconstrucción
levanta la serie de Udenar, \uni{32691}{kWh} son generación devuelta y
\uni{1013}{kWh} proceden del recorte a cero. Elaboración propia aplicando
la ecuación~\eqref{eq:prep-reconstruccion}.}


\begin{detallebox}
Los inversores intervienen en dos capas distintas. La generación $G$ que
el modelo negocia proviene de \emph{un} inversor designado por
institución, que en Udenar es el del proyecto. La reconstrucción de la
ecuación~\eqref{eq:prep-reconstruccion} suma, en cambio, \emph{todos} los
que inyectan detrás del medidor, que en Udenar son tres. Las dos capas
responden a preguntas distintas: cuál generación se negocia, y cuánta
restó el medidor.
\end{detallebox}
\end{guardado}

\begin{guardado}
\subsection{Lo que cuesta reconstruir}
\label{sub:prep-costo}

El recorte a cero de la ecuación~\eqref{eq:prep-reconstruccion} no es
gratuito, y lo que cuesta se puede medir. Actúa cuando la suma de
inversores se queda \emph{corta} frente a la generación que el medidor
descontó, porque solo entonces $D_{\text{neto}} + \sum G$ sigue siendo
negativo; si la suma sobrara, el resultado saldría más positivo, no menos.
La causa es de cobertura: antes del 3 de septiembre de \num{2025} el
inversor del proyecto no registraba, y la suma devuelve al medidor menos
de lo que este había restado.

Lo que eso elimina de la serie de demanda queda acotado. Son las
\num{353} horas de Udenar de la tabla anterior, el \pct{5.7} del
horizonte, todas en la ventana solar y con máximo a las 13 horas, que
suman \uni{1013.3}{kWh}, es decir, el \pct{2.29} de la demanda de Udenar y
el \pct{0.32} de la comunitaria. La distorsión tiene además signo
conocido, porque subestima la demanda de Udenar justo en las horas de
sol.

\begin{lecturabox}
Que el signo sea conocido importa más que la magnitud. Subestimar la
demanda de Udenar en horas solares reduce su déficit y, con él, la energía
que Udenar podría comprar en esas horas. El efecto sobre el mercado entre
pares es, por tanto, restrictivo: hace el mercado más pequeño, no más
grande. Es la primera de una serie de decisiones que apuntan todas en la
misma dirección, y que la subsección~\ref{sub:prep-sesgo} reúne.
\end{lecturabox}
\end{guardado}

\subsection{El umbral de los atípicos}
\label{sub:prep-umbral}

La última etapa que altera el valor del dato es la limpieza, que primero
retira los valores imposibles y después rellena lo que falta. Actúa sobre
las 10 series de medida, que son la demanda y la generación de cada
institución, aunque en la práctica solo sobre la demanda: las cinco de
generación atraviesan la etapa sin que se les toque una sola hora.

Todo se hace a traves del siguiente criterio, un valor se retira cuando supera el umbral dado por:

\begin{equation}
 u = \max\bigl(Q_{75} + 5\,\mathrm{IQR},\; 1{,}2 \cdot P_{99{,}5}\bigr),
 \label{eq:prep-umbral}
\end{equation}

que combina dos criterios. La tensión que resuelve es la siguiente: una
lectura físicamente imposible hay que retirarla, pero un criterio
demasiado estricto retira consumo real, y mirando el valor aislado no hay
forma de distinguirlos.

El primer criterio mide cuánto se dispersa la mitad central de las
lecturas y pone el corte cinco veces esa dispersión por encima de ella.
Sigue la regla de Tukey, la que traza los bigotes de un diagrama de caja,
pero con un multiplicador de 5 en vez del 1,5 clásico, deliberadamente
permisivo para que solo retire lo inverosímil. Funciona bien mientras la
serie tenga recorrido.

Donde falla es en las cargas muy planas. Si la mitad central de las
lecturas cabe en un margen estrecho, el corte queda cerca de esa base y
cualquier pico operativo legítimo lo supera. CESMAG es ese caso: bajo M1 su
mitad central mide poco más de un kilovatio y el corte cae en
\uni{10.6}{kW}, un valor que la institución rebasa con frecuencia.

El segundo criterio existe para eso. Fija un valor por debajo del cual el
corte no puede bajar, y lo sitúa un \pct{20} por encima del percentil 99,5,
es decir, por encima de lo que la propia serie alcanza operando con
normalidad. Como el umbral es el mayor de los dos, se aplica siempre el
\emph{más permisivo}, y el segundo solo entra cuando el primero se ha vuelto
demasiado estricto.

La Figura~\ref{fig:prep-anatomia} construye los dos candidatos sobre las
diez series de demanda, de modo que la multiplicación por cinco se cuenta
con el ojo en vez de leerse en la fórmula.

\begin{figure}[H]
 \centering
 \includegraphics[width=0.95\textwidth]{figuras/f3_04b_anatomia_umbral.png}
 \caption[La anatomía del umbral]{Cómo se construye cada uno de los dos
 candidatos, en las diez series de demanda y en las dos fronteras. La caja
 oscura es la mitad central de las lecturas; sobre el tercer cuartil se
 apilan cinco copias exactas de esa misma caja, y donde aterriza la quinta
 está el primer criterio. A su derecha, la línea de puntos sube hasta el
 percentil 99,5 y el bloque ámbar es su prolongación del \pct{20}, cuyo
 techo es el segundo criterio. El que fija el umbral va con trazo lleno y
 el descartado, discontinuo.}
 \label{fig:prep-anatomia}
\end{figure}

\notafig{Cada panel lleva su propia escala vertical, porque los rangos van
de \uni{0.3}{kW} a \uni{24.4}{kW}, de modo que lo que se compara entre
paneles no es una altura sino una relación, esto es, si la cola con su
prolongación queda por encima o por debajo de la quinta copia. CESMAG bajo
M3 es el caso extremo: su pila apenas despega del suelo del panel y el
umbral queda cuatro veces más arriba, porque su cola cae a 23,2 rangos del
cuartil. El Hospital, con el rango más estrecho de todo el estudio,
\uni{0.318}{kW}, no necesita el segundo criterio en ninguna de las dos
fronteras. Los valores de cada panel están en la
Tabla~\ref{tab:prep-umbral}. Elaboración propia desde el caché de estados
intermedios.}

Lo que está en juego se mide en CESMAG. Sin ese segundo criterio, su
demanda perdería \num{451} horas bajo M1, que concentran el \pct{20.7} de
su energía, y \num{927} bajo M3, que concentran el \pct{52.5}. Con él tiene
cero atípicos bajo M1 y dos bajo M3.

El segundo criterio manda en cuatro de las 20 series, todas de demanda:
Mariana y CESMAG en las dos fronteras. En las 16 restantes basta el
primero, y entre ellas están las diez de generación, donde además no se
retira una sola hora. La Tabla~\ref{tab:prep-umbral} da la anatomía
completa de las series de demanda.

\begin{table}[H]
 \centering
 \caption{Los dos candidatos en cada serie de demanda, con el que se
 aplica en negrita. La columna de la cola dice a cuántos rangos
 intercuartílicos del tercer cuartil cae el percentil 99,5, y es la que
 explica el reparto: cuanto más lejos queda la cola del cuerpo, más se
 queda corto el primer criterio.}
 \label{tab:prep-umbral}
 \footnotesize
 \setlength{\tabcolsep}{5pt}
 \begin{tabular}{@{}l r r r r r@{}}
 \toprule
 \textbf{Institución} & \textbf{Rango} & \textbf{Cola} &
 \textbf{Primer} & \textbf{Segundo} & \textbf{Horas} \\
 & \textbf{(kW)} & \textbf{(rangos)} & \textbf{criterio (kW)} &
 \textbf{criterio (kW)} & \textbf{retiradas} \\
 \midrule
 \multicolumn{6}{@{}l}{\itshape Frontera M1} \\
 Udenar & \num{4,91} & \num{3,4} & \textbf{\num{33,4}} & \num{31,0} & \num{0} \\
 Mariana & \num{4,50} & \num{4,1} & \num{33,5} & \textbf{\num{35,1}} & \num{4} \\
 UCC & \num{24,35} & \num{1,3} & \textbf{\num{154,8}} & \num{78,5} & \num{0} \\
 HUDN & \num{1,40} & \num{1,2} & \textbf{\num{16,8}} & \num{13,7} & \num{0} \\
 CESMAG & \num{1,27} & \num{8,2} & \num{10,6} & \textbf{\num{17,7}} & \num{0} \\
 \addlinespace
 \multicolumn{6}{@{}l}{\itshape Frontera M3} \\
 Udenar & \num{1,54} & \num{2,2} & \textbf{\num{9,5}} & \num{6,3} & \num{0} \\
 Mariana & \num{0,97} & \num{5,0} & \num{7,6} & \textbf{\num{9,1}} & \num{5} \\
 UCC & \num{3,62} & \num{1,2} & \textbf{\num{23,0}} & \num{11,1} & \num{0} \\
 HUDN & \num{0,32} & \num{1,5} & \textbf{\num{3,7}} & \num{3,0} & \num{0} \\
 CESMAG & \num{0,56} & \num{23,2} & \num{4,5} & \textbf{\num{17,6}} & \num{2} \\
 \bottomrule
 \end{tabular}
\end{table}

Que la mitad central sea estrecha no lo decide por sí solo, y la tabla lo
enseña en dos filas contiguas: bajo M1 el HUDN y CESMAG tienen rangos casi
iguales y el umbral de cada uno lo fija un criterio distinto. Lo que los
separa es la columna de la cola, \num{1,2} rangos frente a \num{8,2}. Eso
es lo que obliga al segundo criterio, una base estable con picos altos y
poco frecuentes, y no la estrechez del cuerpo.

El criterio retira \num{11} horas en las 20 series. Todas son de demanda y
solo de dos instituciones: Mariana, con cuatro bajo M1 y cinco bajo M3, y
CESMAG, con dos bajo M3.

Queda una objeción por responder, y es que el multiplicador de 5 no
procede de ninguna referencia sino de la inspección de estas series. La
respuesta no es defender el valor sino medir cuánto depende de él el
resultado. La Tabla~\ref{tab:prep-sensibilidad} barre el multiplicador
entre el clásico y el doble del elegido, con los dos criterios y con el
primero solo.

\begin{table}[H]
 \centering
 \caption{Energía que el criterio retira, en porcentaje de la demanda de
 las 10 series, al variar el multiplicador del primer criterio.}
 \label{tab:prep-sensibilidad}
 \small
 \begin{tabular}{@{}c r r@{}}
 \toprule
 \textbf{Multiplicador} & \textbf{Con los dos} & \textbf{Solo con el} \\
 & \textbf{criterios} & \textbf{primero} \\
 \midrule
 1,5 & \pct{0.083} & \pct{13.369} \\
 3 & \pct{0.083} & \pct{6.563} \\
 \textbf{5} & \textbf{\pct{0.062}} & \pct{3.650} \\
 7 & \pct{0.021} & \pct{2.356} \\
 \bottomrule
 \end{tabular}
\end{table}

La columna de la izquierda contesta la objeción: entre el multiplicador
clásico y el elegido hay cuatro horas de diferencia sobre las \num{122880}
horas de serie del estudio, y la energía afectada no llega en ningún caso
al \pct{0.1} de la demanda. El valor del multiplicador no gobierna nada.

La columna de la derecha explica por qué. Sin el segundo criterio, el
multiplicador sí gobierna, y de forma brutal: el clásico retiraría el
\pct{13.4} de la demanda y ni siquiera un 7 bajaría del \pct{2.4}. El
segundo criterio no cubre un fallo del primero, sino que hace que el
primero deje de importar, y es esa insensibilidad, y no la elección del
valor, lo que sostiene el criterio.

La Figura~\ref{fig:prep-umbral} muestra a la izquierda qué le pasa a una de
esas horas y a la derecha por qué el umbral está donde está en cada serie.

\begin{figure}[H]
 \centering
 \includegraphics[width=0.95\textwidth]{figuras/f3_05_umbral_atipicos_m1.png}
 \caption[El umbral de atípicos y lo que retira]{A la izquierda, la hora
 que el umbral retira con mayor holgura entre las que quedan aisladas, la
 de la Universidad Mariana el 1 de octubre de 2025, con el valor que entra
 y el que queda tras la cascada. A la derecha, el criterio en las diez
 series, con el eje en múltiplos del percentil 99,5 de cada una: la barra
 es la mitad central de las lecturas, la línea llega hasta el máximo
 observado, el círculo marca el corte del primer criterio y el rombo el
 umbral que se aplica.}
 \label{fig:prep-umbral}
\end{figure}

\notafig{La hora del caso entra con \uni{39.9}{kW}, excede en un
\pct{13.7} el umbral de \uni{35.1}{kW} y sale con \uni{28.2}{kW},
interpolada entre las lecturas de las 8 y las 10 horas. En el panel
derecho, el segundo criterio es la vertical en 1,2, común a las diez filas,
de modo que manda justamente donde el círculo cae a su izquierda. La
distancia entre la masa de la serie y su umbral es lo que acredita el
criterio, y solo en cuatro de las 20 series hace falta el segundo para
conservarla. Elaboración propia desde el caché de estados intermedios.}

\subsection{Tratamiento de los huecos}
\label{sub:prep-huecos}

Lo que queda como faltante, sea porque el equipo no registró o porque el
umbral acaba de retirarlo, pasa por una cascada de cuatro tratamientos, de
los que solo tres llegan a actuar. La cascada no clasifica los huecos por
su longitud: se aplica hora a hora dentro de cada uno, de modo que un
mismo hueco puede recibir varios.

Las tres primeras horas de cualquier hueco se interpolan linealmente en el
tiempo, sobre la recta que une la última hora observada antes del hueco
con la primera de después. Si el hueco dura tres horas o menos, ahí
termina. Si es más largo, el tramo interpolado avanza solo esas tres horas
y el resto pasa al tratamiento siguiente.

Ese resto se rellena repitiendo un valor, primero hacia adelante hasta 24
horas y después hacia atrás otras 24. Hacia adelante, el valor que se
repite no es el que cabría esperar: se propaga la tercera hora
interpolada, que es un valor fabricado, y no la última lectura del equipo.
Hacia atrás sí se propaga una lectura real, la primera que reabre la
serie. En los dos casos se repite un único número, de modo que el tramo
queda plano en vez de recuperar el ciclo diario.

Sumados los tres tramos, tres horas más 24 más otras 24, la cascada cierra
cualquier hueco de hasta 51 horas. El cuarto tratamiento imputaría con
cero lo que sobreviviera, y por eso no llega a actuar en este horizonte:
el hueco más largo mide 49 horas, en la demanda de la UCC bajo M3.

La limpieza trata en total \num{808} horas repartidas entre las 20 series,
que son las \num{11} que el umbral retira más las \num{797} que ya
faltaban en el dato. De ellas, \num{233} quedan interpoladas y \num{575},
el \pct{71.2}, rellenas repitiendo un valor.

Que ese relleno sea plano tiene una consecuencia que hay que acotar, y es
que introduce mesetas donde el capítulo va a leer después la coincidencia
entre el pico solar y el valle de demanda. La cota es doble. La primera es
que \emph{la generación no se rellena ni una sola hora} en ninguna de las
dos fronteras, de modo que el pico solar nunca se fabrica: lo que la
limpieza toca es siempre demanda. La segunda es el tamaño de lo tocado.
Sobre las \num{30720} horas de serie de cada frontera, la limpieza rellena
\num{234} bajo M1 y \num{574} bajo M3, y de ellas caen en la franja diurna
\num{121} y \num{300}, es decir, el \pct{0.39} y el \pct{0.98} de las
horas de serie.

La Figura~\ref{fig:prep-huecos} recorre los tres tratamientos sobre huecos
reales de longitud creciente.

\begin{figure}[H]
 \centering
 \includegraphics[width=0.95\textwidth]{figuras/f3_06_escalera_huecos.png}
 \caption[La escalera de los huecos]{Tres huecos reales de longitud
 creciente, con el eje horizontal en horas desde el inicio de cada uno, que
 es la magnitud de la que depende el tratamiento, y cada panel con su
 frontera de medición, porque los tres no salen de la misma. Abajo, cuántas
 horas trató la limpieza en cada institución y en cada frontera, y por qué
 mecanismo.}
 \label{fig:prep-huecos}
\end{figure}

\notafig{El hueco de tres horas lo cierra entero la interpolación. El de 13
sale con sus tres primeras horas interpoladas y las diez restantes
arrastradas, y la meseta no se queda en la última lectura observada,
\uni{9.9}{kW}, sino en la tercera hora interpolada, \uni{10.1}{kW}, es
decir, en un valor que ninguna hora observada tuvo. El de 49, el más largo
de todo el estudio, añade el arrastre hacia atrás desde la lectura que
reabre la serie. Los ceros de M1 en Udenar y en la Universidad Cooperativa
de Colombia no significan que no tuvieran cortes: su medidor es neto y el
recorte de la reconstrucción cerró esos huecos antes de que la limpieza los
viera, y bajo M3 esas mismas dos instituciones acumulan \num{97} y
\num{126} horas. El arrastre es además el \pct{71.2} de todo lo que la
limpieza trata, de modo que la mayor parte de lo rellenado es una meseta y
no una rampa. Elaboración propia desde el caché de estados intermedios.}

\begin{trampabox}
En el proyecto conviven tres criterios de relleno distintos:

\begin{enumerate}[leftmargin=*, itemsep=0.2em, topsep=0.3em]
 \item Las 10 series de medida interpolan las tres primeras horas de cada
 hueco, arrastran el resto hasta 24 horas por cada lado y solo entonces
 recurrirían al cero.
 \item La serie horaria de precios de bolsa, que construye el
 Capítulo~\ref{sec:precios}, interpola las seis primeras horas de cada
 hueco y completa el resto con la mediana de la propia serie, sin marcado
 de atípicos y sin arrastre.
 \item El caso de estudio prospectivo interpola también las seis primeras
 y completa el resto con cero.
\end{enumerate}

Las diferencias son tres y solo una está justificada. Que el residuo de
los precios sea la mediana y no cero se sostiene, porque un consumo de
cero kilovatios es una lectura legítima pero un precio de cero no describe
ninguna hora del mercado mayorista. El límite de interpolación, seis horas
en vez de tres, y la ausencia de arrastre en los dos últimos criterios no
tienen justificación registrada: son heterogeneidad heredada de la
construcción incremental del repositorio, y no elecciones fundamentadas.
\end{trampabox}

\subsection{Ensamblaje y verificación de las matrices}
\label{sub:prep-matrices}

Las 10 series limpias se apilan en dos matrices de dimensión
$5 \times \num{6144}$ y se someten a una verificación de contrato: si
alguna celda de $D$ o de $G$ resulta negativa, la ejecución se detiene con
un error explícito en lugar de continuar con datos inconsistentes. La
no negatividad no es aquí una aspiración sino una propiedad garantizada
por construcción y comprobada al final. La
Figura~\ref{fig:prep-matrices} despliega las dos matrices como mapas de
hora contra día, que es la forma en que sus estructuras se vuelven
visibles de una vez.

\begin{figure}[H]
 \centering
 \includegraphics[width=0.95\textwidth]{figuras/f3_07_matrices_m1.png}
 \caption[Las matrices que consume el modelo]{Las dos matrices
 resultantes, dibujadas como mapas de hora del día contra día del
 horizonte, en la frontera M1. La generación traza una banda diurna
 nítida que sigue el arco solar; la demanda muestra la semana laboral
 en franjas verticales regulares.}
 \label{fig:prep-matrices}
\end{figure}

\begin{lecturabox}
Esta figura funciona como comprobación final del capítulo. Si hubiera un
error de zona horaria, la banda solar aparecería desplazada respecto del
mediodía; si la alineación temporal entre medidores e inversores fallara,
la banda se difuminaría; si la agregación horaria estuviera mal, las
franjas de la semana laboral no se distinguirían. Que las tres estructuras
aparezcan donde deben es lo que autoriza a seguir adelante.
\end{lecturabox}

\subsection{La zona horaria del índice}
\label{sub:prep-zona}

La última etapa no altera ningún valor: le pone zona horaria al índice.
Aun así encierra una suposición de la que dependen dos argumentos
anteriores.

Las marcas de tiempo llegan sin zona, de modo que hay que decidir cuál es
la suya. El pipeline las declara hora local de Colombia, es decir, las
interpreta como ya escritas en ese huso en vez de convertirlas desde otro.
La suposición es que cada equipo escribe la hora del sitio donde está
instalado, y se comprueba sin salir del dato: el máximo de generación de
las cinco instituciones cae entre las 11 y las 12, que es el mediodía solar
de Pasto. Si las marcas vinieran en tiempo universal, con Colombia cinco
horas por detrás, ese máximo aparecería entre las 16 y las 17.

Los dos argumentos que dependen de esto son el de la
subsección~\ref{sub:prep-negativa}, que sitúa las lecturas negativas en las
horas de sol, y el de la subsección~\ref{sub:prep-matrices}, que lee la
banda solar de la matriz de generación en torno al mediodía. Los dos se
caerían si la serie estuviera desplazada cinco horas.

Colombia no aplica horario de verano, de modo que en este conjunto no hay
horas que no existan ni horas que ocurran dos veces. Las dos salvaguardas
que el pipeline lleva para esos casos, una que adelanta la hora inexistente
y otra que resuelve la ambigua, nunca llegan a actuar.

\subsection{Cinco ritmos distintos}
\label{sub:prep-perfiles}

Con las matrices ya construidas se puede mirar lo que caracteriza a esta
comunidad y hace posible el intercambio: las cinco instituciones comparten
municipio, clima y recurso solar, pero no comparten rutina. La
Figura~\ref{fig:prep-perfiles} pone los cinco perfiles uno junto a otro,
con el de la comunidad completa como sexto panel.

\begin{figure}[H]
 \centering
 \includegraphics[width=0.95\textwidth]{figuras/f3_08_perfiles_instituciones_m1.png}
 \caption[Perfil medio de cada institución]{Perfil medio de demanda y
 generación por institución en la frontera M1; el porcentaje del
 título es la razón entre generación y demanda de cada una. La curva
 verde sombreada es común en forma a las cinco (el sol es el mismo),
 mientras que las curvas de demanda difieren en nivel y en forma.}
 \label{fig:prep-perfiles}
\end{figure}

\notafig{El HUDN, que es un hospital, mantiene una demanda casi plana
las 24 horas, como corresponde a una instalación que no cierra.
Udenar dibuja el doble pico característico de un campus universitario,
con el valle del almuerzo entre ambos. La UCC opera al nivel más alto del
conjunto en la frontera M1 y concentra su consumo en la jornada diurna.
Elaboración propia desde el caché de estados intermedios.}

Esa disparidad de formas es exactamente la condición que hace que unas
instituciones tengan excedente cuando otras tienen déficit, y sin ella no
habría nada que intercambiar.

\subsection{El ritmo semanal y el anual}
\label{sub:prep-ritmos}

Hay una segunda fuente de disparidad, esta vez temporal y común a las
cinco: el calendario. La Figura~\ref{fig:prep-ritmos} cruza el perfil
diario con el corte semanal y con el recorrido mensual.

\begin{figure}[H]
 \centering
 \includegraphics[width=0.95\textwidth]{figuras/f3_09_ritmos_m1.png}
 \caption[Ritmo semanal y recorrido anual]{A la izquierda, día hábil
 frente a fin de semana: la demanda cae de forma apreciable mientras la
 generación permanece igual. A la derecha, el recorrido mensual de las
 dos magnitudes a lo largo del horizonte.}
 \label{fig:prep-ritmos}
\end{figure}

\notafig{El corte semanal esconde una asimetría que reaparece en los
resultados: el fin de semana la demanda cae pero la generación no, porque
el sol no distingue entre sábado y martes. La consecuencia es que el
excedente disponible para intercambio es mayor precisamente cuando hay
menos gente en los edificios, es decir, cuando la comunidad tiene menos
capacidad de absorberlo internamente. Elaboración propia desde el caché
de estados intermedios.}

\subsection{Las decisiones que sesgan en contra}
\label{sub:prep-sesgo}

En las seis etapas se decidió algo y el resultado sale favorable, de modo
que la objeción es inmediata: cuánto del resultado son las decisiones. La
respuesta honesta no es que las decisiones sean neutrales, sino que lo que
importa es hacia dónde apuntan.

Seis apuntan en la misma dirección, la de recortar el resultado. Cuatro son
del pipeline y dos del modelo:

\begin{enumerate}[leftmargin=*, itemsep=0.35em]
 \item \textbf{El recorte de la reconstrucción} elimina
 \uni{1013.3}{kWh} de la demanda de Udenar en horas solares, según la
 Tabla~\ref{tab:prep-reconstruccion}. Menos demanda en horas de sol
 significa menos déficit que cubrir y, por tanto, un mercado más
 pequeño.

 \item \textbf{Un solo medidor y un solo inversor por institución.} El
 pipeline no suma todos los equipos disponibles sino que designa uno de
 cada clase. En Udenar, que aporta tres inversores y expone uno, la
 generación que ve el mercado queda por debajo de la instalada; en las
 otras cuatro el equipo designado es el único, de modo que coinciden.
 En el agregado comunitario, por tanto, se expone menos de lo que hay.

 \item \textbf{La ausencia de gestión de la demanda.} El modelo del
 Capítulo~\ref{sec:modelo} corre con la demanda tal como se midió, sin
 programa de desplazamiento de carga. Toda ganancia que un despliegue
 real obtendría moviendo consumo hacia las horas de sol queda fuera del
 resultado.

 \item \textbf{El descarte conservador de las horas numéricamente
 difíciles.} Cuando el método numérico que resuelve el equilibrio
 de una hora no converge, esa hora se descarta sin imputarle beneficio
 al mercado entre pares, en vez de estimarlo.

 \item \textbf{La generación que falta se pone a cero antes de la
 limpieza}, de modo que sus huecos no se interpolan ni se arrastran como
 los de la demanda. En la noche eso es correcto, porque el equipo no
 produce, pero \num{285} horas de sol del horizonte quedan con generación
 nula cuando lo que hubo fue una avería de registro. Menos generación es
 menos excedente que ofrecer.

 \item \textbf{La demanda sin lectura de medidor también se pone a cero},
 en la reconstrucción y no en la limpieza, según la
 subsección~\ref{sub:prep-reconstruccion}. Son \num{330} horas bajo M1 que
 entran al modelo como consumo nulo, y menos demanda es menos déficit que
 cubrir.
\end{enumerate}

La séptima apunta al revés. Al estimar la hora incompleta con la media de
sus muestras presentes se le atribuye al tramo no observado el
comportamiento del observado, lo que da entre \pct{0.26} y \pct{0.55} más
energía que truncarlo. Afecta a \num{835} horas del horizonte, es decir, al
\pct{2.8} de las horas con dato, en las que se observa de media el
\pct{85.9} de la hora.

No se elige por su signo: truncar queda desmentido por el contador del
propio equipo y va en contra del criterio del regulador, según la
subsección~\ref{sub:prep-incompleta}. El supuesto se comprobó además
aplicando la ventana observada de cada hora incompleta a horas completas
comparables de la misma institución, la misma hora del día y el mismo mes:
el sesgo resulta indistinguible de cero en las cinco, y suma el
\pct{0.0007} de la demanda del horizonte.

Seis de las siete recortan el resultado y la restante lo amplía en una
proporción tres órdenes de magnitud menor que el margen que separa a los
mecanismos. La
consecuencia metodológica es que las cifras del Capítulo~\ref{sec:res}
deben leerse como una \textbf{cota inferior} de lo que este mercado podría
rendir, y no como una estimación centrada. Es una posición deliberadamente incómoda, porque renuncia a parte del
margen, y es la que permite sostener el resultado ante una lectura
hostil.

\subsection{Declaración de honestidad metodológica}
\label{sub:prep-honestidad}

Cuatro elecciones de este capítulo se declaran sin atenuarlas. Las tres
primeras carecen de respaldo formal; la cuarta está adoptada pero todavía
no aplicada.

\begin{enumerate}[leftmargin=*, itemsep=0.35em]
 \item El umbral de valores atípicos
 $\max(Q_{75} + 5\,\mathrm{IQR},\; 1{,}2 \cdot P_{99{,}5})$ se adoptó
 por inspección de las series del proyecto, y no hay referencia de
 literatura que respalde el multiplicador cinco. Lo que sí hay es la
 medida de cuánto depende el resultado de él, según la
 Tabla~\ref{tab:prep-sensibilidad}: entre el multiplicador clásico y el
 elegido la energía retirada varía dos centésimas de punto porcentual, de
 modo que la falta de respaldo bibliográfico no compromete la cifra. Lo
 que sostiene el criterio no es el valor sino esa insensibilidad.

 \item El relleno hacia adelante con límite de 24 horas es una
 heurística de «día operativo» (supone que el patrón del día anterior es
 representativo) cuyo impacto no se evaluó de forma aislada.

 \item No se construyó una matriz de sensibilidad del pipeline completo,
 es decir, variar el umbral y los límites de interpolación y relleno y
 medir el efecto sobre los resultados. Queda como trabajo futuro.
 \item Se adopta como criterio de validez que una hora traiga al menos el
 \pct{75} de sus muestras, y la corrida que este documento reporta es
 anterior a esa adopción, de modo que no lo aplica: las \num{147} horas
 que no lo cumplen entraron estimadas como las demás. El efecto está
 medido y vale el \pct{0.026} de la demanda comunitaria, según la
 subsección~\ref{sub:prep-incompleta}. Incorporarlo exige rehacer la corrida
 canónica completa, de modo que queda pendiente para la próxima, y hasta
 entonces la cifra de arriba es la cota de lo que cambiaría.
\end{enumerate}

Como cota global, el impacto neto de estas elecciones sobre el volumen de
energía transada se estima en $\pm$\,\pct{0.3}. El estatuto de esa cifra
necesita una precisión, porque es fácil leerla como un resultado y no lo
es: no procede de la matriz de sensibilidad que el
punto~3 declara no construida, sino de un juicio apoyado en dos evidencias indirectas: la cobertura de
adquisición por equipo, que supera el \pct{97} en los inversores
principales, y la magnitud del recorte de la reconstrucción, que
representa el \pct{0.32} de la demanda comunitaria. Es una estimación
declarada y no un valor calculado, y así debe leerse.

Con todo lo anterior, el dato queda listo para el modelo: dos matrices de
\num{6144} horas, no negativas por construcción, con la energía que cada
institución consume y la que produce separadas de forma explícita, y con
el costo de cada intervención cuantificado. Queda, sin embargo, una decisión
que este capítulo ha mencionado sin desarrollar, la de qué medidor se lee
en cada institución, y de ella depende buena parte de lo que sigue: es el
objeto del capítulo siguiente.
